%% kapa: exam=1 round=25 year=2014 period=1 form=A subject=civil_law subject_ko=민법 seq=37 number=37
%% source: Q-Net 감정평가사 기출문제 (cst00309) artlSeq=5101207 file=extracted/1008/5101207/2014년 제25회 감정평가사 1차시험 기출문제/1교시(A형).hwp
%% conversion: kordoc 4.24.51 -> tools/gen1.py (auto); review=proofread-sonnet; flags=none
\begin{question}{37}
지역권에 관한 설명으로 옳지 않은 것은?
\begin{choices}
  \item 승역지에 수개의 용수지역권이 설정된 때에는 후순위의 지역권자는 선순위의 지역권자의 용수(用水)를 방해하지 못한다.
  \item 요역지가 수인의 공유인 경우에 그 공유자 1인에 의한 지역권 소멸시효의 중단은 다른 공유자를 위하여 효력이 있다.
  \item 지역권은 계속되고 표현된 것에 한하여 시효취득할 수 있다.
  \item 지역권설정계약에 의한 지역권 취득은 등기하여야 그 효력이 생긴다.
  \item 지역권은 요역지와 분리하여 다른 권리의 목적으로 할 수 있다.
\end{choices}
\end{question}
\answer{5}
